\section{Metodologia}\label{sec:metodologia}

Este projeto realizará a avaliação de diferentes técnicas de extração de características para reconhecimento de usuários pela impressão digital. Para isso, serão utilizados conjuntos de dados com imagens de impressão digital. A partir dessas imagens serão aplicadas técnicas de extração de características para realização da avaliação deste projeto.

Alguns conjuntos de dados com dados publicamente disponíveis são apresentados a seguir:

\begin{itemize}
    \item \textit{Fingerprint Verification Competition} (FVC): foram realizadas algumas competições organizadas por pesquisadores da \textit{University of Bologna} para avaliação de sistemas de reconhecimento de usuários pela impressão digital. Os conjuntos de dados da parte B dos anos 2000, 2002 e 2004 estão disponíveis para download \footnote{\url{http://bias.csr.unibo.it/fvc2000/download.asp}} 	\footnote{\url{http://bias.csr.unibo.it/fvc2002/download.asp}} 	\footnote{\url{http://bias.csr.unibo.it/fvc2004/download.asp}}.
    
    \item \textit{Neurotechnology}: a empresa Neurotechnology disponibilizou dois conjuntos de imagens de impressão digital, que podem ser acessados pelo site \footnote{\url{https://www.neurotechnology.com/download.html\#databases}}. O primeiro conjunto de dados possui imagens capturadas a partir do leitor \textit{Cross Match Verifier 300} enquanto que o segundo possui imagens obtidas a partir do leitor \textit{DigitalPersona U.are.U 4000}.
\end{itemize}
    
Durante a avaliação, poderá ser avaliado o desempenho preditivo de algoritmos de classificação. Neste caso, os dados disponíveis serão divididos entre treino e teste. Para avaliação do desempenho preditivo serão usadas métricas para avaliação comumente adotadas na literatura, como FMR (\textit{False Match Rate}), FNMR (\textit{False Non-Match Rate}) e acurácia balanceada \cite{Precise2014, Ferlini2021eargate}. FMR mede o percentual de tentativas de impostores incorretamente classificadas como genuínas, enquanto que a FNMR mede o percentual de tentativas genuínas incorretamente classificadas como impostoras. A acurácia balanceada representa um desempenho médio considerando o mesmo peso para tentativas impostores e genuínas.